\documentclass[11pt,a4paper]{article}

% Essential packages
\usepackage[utf8]{inputenc}
\usepackage[british]{babel}
\usepackage{amsmath,amssymb,amsthm}
\usepackage{geometry}
\usepackage{graphicx}
\usepackage{booktabs}
\usepackage{hyperref}
\usepackage{xcolor}
\usepackage{enumitem}

% Page geometry
\geometry{
    a4paper,
    left=2.5cm,
    right=2.5cm,
    top=2.5cm,
    bottom=2.5cm
}

% Hyperref settings
\hypersetup{
    colorlinks=true,
    linkcolor=blue,
    citecolor=blue,
    urlcolor=blue
}

% Title information
\title{\textbf{Error Analysis Theory for MLMC and Laplace Methods in American Basket Option Pricing}}
\author{}
\date{}

\begin{document}

\maketitle

\noindent\textbf{Estimating numerical error in high-dimensional option pricing requires a sophisticated framework combining telescoping sum decomposition, Richardson extrapolation, and primal-dual bounds}---essential tools when ground truth is unknowable. This report provides the mathematical formulations and practical implementation guidance needed for rigorous error quantification in Markovian projection methods with L² polynomial regression.

\vspace{0.3cm}

\noindent The core challenge in American basket option pricing is that \textbf{neither MLMC nor Laplace approximation provides exact solutions}, yet practitioners must quantify computational accuracy. Modern error analysis resolves this through three complementary approaches: the MLMC telescoping sum framework estimates bias from level differences without knowing truth; Laplace approximation error is controlled through primal-dual bounds; and cross-validation between methods establishes confidence when analytical benchmarks don't exist.

\section{The MLMC telescoping sum framework for error estimation}

The foundational MLMC identity exploits a telescoping structure that enables error estimation without knowing the true value $P$:

\begin{equation}
\mathbb{E}[P_L] = \mathbb{E}[P_0] + \sum_{\ell=1}^{L} \mathbb{E}[P_\ell - P_{\ell-1}]
\end{equation}

\noindent where $P_\ell$ denotes the level-$\ell$ numerical approximation with increasing accuracy as $\ell$ increases. The unbiased MLMC estimator becomes:

\begin{equation}
\hat{Y} = \sum_{\ell=0}^{L} \hat{Y}_\ell = N_0^{-1}\sum_{n=1}^{N_0} P_0^{(0,n)} + \sum_{\ell=1}^{L} \left\{ N_\ell^{-1} \sum_{n=1}^{N_\ell} \left(P_\ell^{(\ell,n)} - P_{\ell-1}^{(\ell,n)}\right) \right\}
\end{equation}

\noindent The critical insight is that \textbf{$P_\ell^{(\ell,n)}$ and $P_{l-1}^{(\ell,n)}$ must use the same underlying stochastic sample} (same Brownian path)---this coupling creates the correlation that drives variance reduction. For the mean square error decomposition:

\begin{equation}
\text{MSE} = \underbrace{\sum_{\ell=0}^{L} N_\ell^{-1} V_\ell}_{\text{Variance (statistical error)}} + \underbrace{(\mathbb{E}[P_L] - \mathbb{E}[P])^2}_{\text{Bias}^2}
\end{equation}

\noindent The key assumptions enabling tractable analysis are \textbf{geometric convergence rates}: bias decays as $|\mathbb{E}[P_\ell - P]| \leq c_1 \cdot 2^{-\alpha\ell}$ with weak convergence rate $\alpha$, while variance decays as $V_\ell \leq c_2 \cdot 2^{-\beta\ell}$ with rate $\beta$.

\section{Estimating bias without knowing truth via Richardson extrapolation}

Since the true value $P$ is unknown, the \textbf{level differences serve as a proxy for bias estimation}. The Richardson extrapolation-like technique proceeds from the assumption that weak error follows $\mathbb{E}[P_\ell - P] = c \cdot 2^{-\alpha\ell} + O(2^{-2\alpha\ell})$. This yields:

\begin{equation}
\mathbb{E}[P - P_L] = \sum_{\ell=L+1}^{\infty} \mathbb{E}[P_\ell - P_{\ell-1}] \approx \frac{\mathbb{E}[P_L - P_{L-1}]}{2^\alpha - 1}
\end{equation}

\noindent The practical bias estimate therefore uses the observable quantity $|m_L| = |\mathbb{E}[P_L - P_{L-1}]|$ scaled by the factor $(2^\alpha - 1)^{-1}$. Giles' implementation extends this to a robust estimator using multiple previous levels:

\begin{equation}
\text{bias estimate} = \max_{\ell \in \{L-2, L-1, L\}} \frac{|m_\ell|}{2^{\alpha(L-\ell)}(2^\alpha - 1)}
\end{equation}

\noindent The weak convergence rate $\alpha$ is estimated empirically via linear regression on $\log_2|m_\ell|$ versus level $\ell$, with a floor of $\alpha \geq 0.5$ for robustness. For Euler-Maruyama discretisation, $\alpha = 1$ is typical; Milstein schemes achieve $\alpha = 2$ for smooth payoffs.

\section{Variance estimation distinguishes ``square of difference'' from ``difference of squares''}

A subtle but critical distinction exists between two variance-related quantities. The \textbf{correct MLMC variance estimator} computes the sample variance of the coupled difference $Y_\ell = P_\ell - P_{l-1}$:

\begin{equation}
\hat{V}_\ell = \frac{1}{N_\ell}\sum_{n=1}^{N_\ell}(Y^{(n)})^2 - \left(\frac{1}{N_\ell}\sum_{n=1}^{N_\ell}Y^{(n)}\right)^2
\end{equation}

\noindent This equals $\mathbb{E}[(P_\ell - P_{\ell-1})^2] - (\mathbb{E}[P_\ell - P_{\ell-1}])^2$---the \textbf{variance of the difference}. The alternative quantity $\mathbb{E}[P_\ell^2] - \mathbb{E}[P_{\ell-1}^2]$ is the \textbf{difference of second moments}, which has fundamentally different meaning.

\vspace{0.2cm}

\noindent The crucial relationship involves the correlation $\rho$ between levels:

\begin{equation}
\text{Var}[P_\ell - P_{\ell-1}] = \text{Var}[P_\ell] + \text{Var}[P_{\ell-1}] - 2\rho\sqrt{\text{Var}[P_\ell]\text{Var}[P_{\ell-1}]}
\end{equation}

\noindent Strong coupling (same Brownian path) makes $\rho \to 1$, so $V_\ell \ll \text{Var}[P_\ell] + \text{Var}[P_{\ell-1}]$. This high correlation is what enables MLMC's variance reduction---the variance reduction factor $\text{VRF}_\ell = (\text{Var}[P_\ell] + \text{Var}[P_{\ell-1}])/V_\ell$ should satisfy $\text{VRF} \gg 1$ (typically \textbf{$>10$} for effective coupling).

\section{L² and L$^\infty$ error norms require pathwise coupling}

The L² error norm in MLMC measures \textbf{strong (pathwise) convergence}:

\begin{equation}
\|P_\ell - P\|_{L^2} = \sqrt{\mathbb{E}[(P_\ell - P)^2]} = \sqrt{\text{Var}[P_\ell - P] + (\mathbb{E}[P_\ell - P])^2}
\end{equation}

\noindent For Lipschitz payoffs $P = f(S)$ with $|f(S_1) - f(S_2)| \leq K\|S_1 - S_2\|$, the variance bound becomes:

\begin{equation}
V[P_\ell - P_{\ell-1}] \leq K^2 \mathbb{E}[\|S_\ell - S_{\ell-1}\|^2]
\end{equation}

\noindent This connects variance decay rate $\beta$ directly to strong convergence: \textbf{Euler-Maruyama achieves $\beta = 1$} (strong error $O(h^{1/2})$), while \textbf{Milstein achieves $\beta = 2$} (strong error $O(h)$). The complexity theorem states that for MSE $< \varepsilon^2$:

\vspace{0.3cm}

\begin{center}
\begin{tabular}{@{}ll@{}}
\toprule
\textbf{Regime} & \textbf{Complexity} \\
\midrule
$\beta > \gamma$ (variance dominates) & $O(\varepsilon^{-2})$ \\
$\beta = \gamma$ & $O(\varepsilon^{-2}(\log\varepsilon)^2)$ \\
$\beta < \gamma$ (cost dominates) & $O(\varepsilon^{-2-(\gamma-\beta)/\alpha})$ \\
\bottomrule
\end{tabular}
\end{center}

\vspace{0.3cm}

\noindent The $L^\infty$ norm $\|P_\ell - P\|_\infty = \sup_\omega |P_\ell(\omega) - P(\omega)|$ is more stringent but rarely used in MLMC since pathwise supremum estimates are difficult to compute practically.

\section{Weak versus strong error exhibits distinct behaviour in regression problems}

\textbf{Weak error} measures distributional convergence: $|\mathbb{E}[P_\ell] - \mathbb{E}[P]|$---the difference of means. \textbf{Strong error} measures pathwise convergence: $\mathbb{E}[|P_\ell - P|^2]^{1/2}$---requiring the same underlying $\omega$.

\vspace{0.2cm}

\noindent For standard SDE discretisation, weak convergence typically achieves one order higher than strong (Euler: weak $O(h)$, strong $O(h^{1/2})$). However, \textbf{regression problems in Longstaff-Schwartz style algorithms introduce complications}:

\begin{itemize}[leftmargin=*,topsep=5pt,itemsep=3pt]
\item The ``truth'' $P$ involves function approximation, not just discretisation
\item Strong convergence assumptions may not hold without explicit pathwise coupling
\item Regression coefficients introduce additional bias-variance trade-offs
\item Level differences may not decay geometrically if regression basis changes across levels
\end{itemize}

\noindent For L² polynomial regression on volatility surfaces, the practical recommendation is to \textbf{use identical basis functions across all MLMC levels} and verify that $\log_2 V_\ell$ versus $\ell$ maintains linear decay with slope $\approx -\beta$.

\section{Laplace approximation error in Markovian projection}

The Bayer-Häppölä-Tempone (2017) framework addresses American basket options by projecting a $d$-dimensional basket onto a 1D Markovian process via Gyöngy's mimicking theorem. For basket index $I_t = \sum_i w_i S_t^i$, the local volatility of the projected process requires computing:

\begin{equation}
\sigma^2(t,z) = \mathbb{E}[\delta_t^2 | I_t = z]
\end{equation}

\noindent where $\delta_t$ is instantaneous volatility. The \textbf{Laplace approximation transforms this $d$-dimensional integral into evaluation at the saddle point}, achieving \textbf{linear complexity in dimension} rather than exponential.

\vspace{0.2cm}

\noindent The Laplace error structure is well-characterised asymptotically: \textbf{leading-order error is $O(t)$ for small $t$} in the short-time regime. For high-dimensional settings where dimension $p = o(n^{1/3})$, the error scales as $O(p^3/n)$ under mild regularity assumptions. The paper does not provide explicit error bounds for the Laplace component; instead, \textbf{total error is controlled through primal-dual bounds}.

\section{Primal-dual bounds quantify total methodology error}

Following Rogers (2002), the duality framework provides rigorous bounds without requiring ground truth:

\begin{equation}
V_0(X_0) = \sup_\tau \mathbb{E}[h_\tau(X_\tau)] = \inf_M \mathbb{E}\left[\max_k (h_k(X_k) - M_k)\right]
\end{equation}

\noindent Any stopping rule $\tau$ yields a \textbf{lower bound} (primal); any martingale $M$ with $M_0 = 0$ yields an \textbf{upper bound} (dual). The practical implementation uses:

\begin{itemize}[leftmargin=*,topsep=5pt,itemsep=3pt]
\item \textbf{Lower bound}: Apply the near-optimal exercise strategy from the projected problem to the full basket---directly simulable
\item \textbf{Upper bound}: Solve only the low-dimensional optimal control problem and construct the dual martingale
\end{itemize}

\noindent \textbf{The bound gap quantifies total error} from all sources: Laplace approximation, Markovian projection, discretisation, and Monte Carlo sampling. For baskets up to $d=50$ assets, the Bayer paper reports \textbf{relative price errors of only a few per cent} with tight bound gaps.

\section{Optimal transport coupling for MLMC variance reduction}

Beyond simple same-random-number coupling, optimal transport provides a principled framework. The Kantorovich formulation seeks:

\begin{equation}
\min_{\pi \in \Pi(\mu, \nu)} \int c(x,y) \, d\pi(x,y)
\end{equation}

\noindent subject to marginal constraints. For MLMC, we want $(X_\ell, X_{\ell-1})$ to have correct marginals but maximal correlation. The \textbf{Giles-Szpruch antithetic coupling} achieves this for multi-dimensional SDEs without Lévy area simulation:

\begin{equation}
Y_\ell = N_\ell^{-1}\sum_i \left[\frac{1}{2}(P(X^f) + P(X^a)) - P(X^c)\right]
\end{equation}

\noindent where $X^f$ (fine) and $X^a$ (antithetic) swap Brownian increments within coarse timesteps. This achieves \textbf{$\beta = 2$ for smooth payoffs} and $\beta = 3/2$ for piecewise smooth payoffs (vanilla calls/puts)---enabling $O(\varepsilon^{-2})$ complexity even with only $O(h^{1/2})$ strong convergence.

\vspace{0.2cm}

\noindent \textbf{Variance reduction verification} requires diagnostic plots:

\vspace{0.3cm}

\begin{center}
\begin{tabular}{@{}ll@{}}
\toprule
\textbf{Diagnostic} & \textbf{Expected Behaviour} \\
\midrule
$\log_2(V_\ell)$ vs $\ell$ & Linear with slope $-\beta$ \\
Correlation $\rho_\ell$ & $> 0.95$ (often $> 0.99$) \\
Variance reduction factor & $> 10$ \\
Kurtosis $\kappa_\ell$ & $< 100$ for reliability \\
\bottomrule
\end{tabular}
\end{center}

\section{Comparing MLMC and Laplace without ground truth}

When neither method provides truth, \textbf{convergence studies combined with Bland-Altman analysis} establish confidence. Richardson extrapolation estimates individual method errors:

\begin{equation}
\text{Error}_h \approx \frac{f_h - f_{2h}}{2^p - 1}
\end{equation}

\noindent for a method of order $p$ with refinement ratio 2. The \textbf{Grid Convergence Index} provides uncertainty quantification:

\begin{equation}
\text{GCI} = 1.25 \cdot |f_h - f_{2h}|/(2^p - 1)
\end{equation}

\noindent For method comparison, the Bland-Altman limits of agreement are essential:

\begin{itemize}[leftmargin=*,topsep=5pt,itemsep=3pt]
\item \textbf{Mean difference (bias)}: $\bar{d} = \frac{1}{n}\sum_i (M_1^{(i)} - M_2^{(i)})$
\item \textbf{95\% Limits of Agreement}: $\bar{d} \pm 1.96 \cdot s_d$
\end{itemize}

\noindent Methods agree if the LOA falls within pre-specified practical tolerances. For volatility surfaces, the \textbf{L² norm over the surface} provides aggregate comparison:

\begin{equation}
\|\sigma_1 - \sigma_2\|_{L^2} = \sqrt{\int\int |\sigma_1(K,T) - \sigma_2(K,T)|^2 \, dK \, dT}
\end{equation}

\noindent with discrete approximation using appropriate quadrature weights. Industry benchmarks suggest \textbf{IV-RMSE $< 1\%$} for liquid strikes/maturities is achievable with high-quality implementations.

\section{Practical implementation guidance for computational physicists}

For American basket option pricing using Markovian projection with L² polynomial regression:

\vspace{0.2cm}

\noindent \textbf{MLMC configuration}: Use antithetic coupling to achieve $\beta = 2$. Start with $L = 4$ levels, refining if bias estimate exceeds $\sqrt{\theta}\varepsilon$ where $\theta \in [0.25, 0.5]$ allocates error budget between bias and variance. Sample allocation follows:

\begin{equation}
N_\ell = 2\varepsilon^{-2}\sqrt{V_\ell/C_\ell} \sum_{\ell'=0}^{L}\sqrt{V_{\ell'}C_{\ell'}}
\end{equation}

\noindent \textbf{Regression basis}: Use \textbf{5-20 weighted Laguerre or Hermite polynomials}---same basis across all levels for consistent strong convergence. Apply regression only to in-the-money paths following Longstaff-Schwartz. Verify fit quality through out-of-sample testing on independent paths.

\vspace{0.2cm}

\noindent \textbf{Error reporting protocol}: Report the Richardson-extrapolated value with GCI uncertainty bands for each method. Compare methods using Bland-Altman plots with bootstrap confidence intervals. Present RMSE stratified by moneyness (ITM/ATM/OTM) and maturity (short/medium/long). State all convergence orders---theoretical and observed.

\vspace{0.2cm}

\noindent \textbf{Diagnostic verification}: Confirm $\log_2 V_\ell$ versus $\ell$ maintains linear slope; correlation $\rho_\ell > 0.95$; kurtosis $\kappa_\ell < 100$; and consistency check $(a - b + c)/\sigma < 1$ where $a, b, c$ are level estimates with combined standard deviation $\sigma$.

\section{Conclusion}

Error analysis for MLMC and Laplace methods in American basket option pricing rests on \textbf{three complementary pillars}: the telescoping sum framework that enables bias estimation from observable level differences rather than unknowable truth; the primal-dual bound methodology that quantifies total Laplace/projection error without explicit component decomposition; and the cross-validation toolkit (Richardson extrapolation, Bland-Altman analysis, convergence studies) that establishes confidence when comparing methods without analytical benchmarks.

\vspace{0.2cm}

\noindent The key implementation insight is that \textbf{variance of level differences}---not difference of variances---governs MLMC efficiency, with strong coupling creating the correlation that enables $O(\varepsilon^{-2})$ complexity. For Markovian projection, the bound gap serves as the practical error metric, implicitly capturing Laplace approximation error alongside all other numerical approximations. When ground truth is absent, self-convergence studies with Richardson extrapolation provide the most rigorous error estimates, while Bland-Altman analysis quantifies method agreement without privileging either as reference.

\end{document}
